% GENERATED FILE. Edit canonical lessons, then run npm run generate:books.
% canonical-source-hash: fnv1a64:63e8004edf21a369
% canonical-lessons: IT-C403-pubblicare, IT-C403-commentare, IT-C403-inoltrare, IT-C403-email, IT-C403-casella

\chapter{Sui social}
\label{ch:it-sui-social}
\hlchaptermodality{403}

\begin{chapteropening}
\textbf{By the end of this chapter:} \emph{I can talk about posting online: publishing, commenting, forwarding a message, email and my inbox.}

Complete the last lesson of chapter 403: I can talk about posting online: publishing, commenting, forwarding a message, email and my inbox.
\end{chapteropening}

\section[pubblicare]{pubblicare — to publish, to post}

\label{lesson:IT-C403-pubblicare}



\begin{quote}
Before the new one: say the Italian for to browse, then the Italian for the link.
\end{quote}

\subsection*{You'll want to know: pubblicare}
\textbf{pubblicare} — \textquotedblleft{}to publish, to post\textquotedblright{}.

\begin{quote}
Se pubblichi la foto, tutti la vedono.
\end{quote}

\begin{quote}
\emph{If you post the photo, everyone sees it.}
\end{quote}

For books too: \emph{pubblicare un romanzo}.

\begin{cousinweb}
\emph{pubblicare} is a learned borrowing from Latin \emph{pūblicāre} \textquotedblleft{}to make public\textquotedblright{} (Lewis \& Short), from \emph{pūblicus} \textquotedblleft{}of the people, public\textquotedblright{}.
\end{cousinweb}

\begin{grammarlens}[title={What you've built}]
One more word for telling, explaining and describing at B1.
\end{grammarlens}

\subsection*{Your turn}
\begin{itemize}
\raggedright
  \item \emph{Say it:} \emph{pubblicare}
  \item \emph{Say it:} \emph{pubblicare}, once more
  \item \emph{Say it:} a sentence of your own with \emph{pubblicare}
  \item \emph{Recall:} say the Italian for to browse, then the Italian for the link, then say \emph{pubblicare} again
\end{itemize}

\begin{tcolorbox}[breakable,colback=teal!4,colframe=teal!35!black,title={Before you move on}]
What does pubblicare mean? (\textquotedblleft{}To publish, to post\textquotedblright{}.) Say it once more.
\end{tcolorbox}
\section[commentare]{commentare — to comment (on)}

\label{lesson:IT-C403-commentare}



\begin{quote}
Before the new one: say the Italian for the link, then the Italian for to publish.
\end{quote}

\subsection*{You'll want to know: commentare}
\textbf{commentare} — \textquotedblleft{}to comment (on)\textquotedblright{}.

\begin{quote}
Molti hanno commentato la sua foto.
\end{quote}

\begin{quote}
\emph{Many people commented on her photo.}
\end{quote}

\emph{Commentare una partita} — to commentate a match.

\begin{cousinweb}
\emph{commentare} is a learned borrowing from Latin \emph{commentārī} \textquotedblleft{}to think over, to make notes on\textquotedblright{} (EtymDB), a verb formed on \emph{comminīscī} \textquotedblleft{}to devise, to recall\textquotedblright{}.
\end{cousinweb}

\begin{grammarlens}[title={What you've built}]
One more word for telling, explaining and describing at B1.
\end{grammarlens}

\subsection*{Your turn}
\begin{itemize}
\raggedright
  \item \emph{Say it:} \emph{commentare}
  \item \emph{Say it:} \emph{commentare}, once more
  \item \emph{Say it:} a sentence of your own with \emph{commentare}
  \item \emph{Recall:} say the Italian for the link, then the Italian for to publish, then say \emph{commentare} again
\end{itemize}

\begin{tcolorbox}[breakable,colback=teal!4,colframe=teal!35!black,title={Before you move on}]
What does commentare mean? (\textquotedblleft{}To comment (on)\textquotedblright{}.) Say it once more.
\end{tcolorbox}
\section[inoltrare]{inoltrare — to forward (a message); to submit}

\label{lesson:IT-C403-inoltrare}



\begin{quote}
Before the new one: say the Italian for to publish, then the Italian for to comment.
\end{quote}

\subsection*{You'll want to know: inoltrare}
\textbf{inoltrare} — \textquotedblleft{}to forward (a message); to submit\textquotedblright{}.

\begin{quote}
Ti inoltro il messaggio della professoressa.
\end{quote}

\begin{quote}
\emph{I will forward you the teacher's message.}
\end{quote}

In offices, \emph{inoltrare una domanda} is to submit an application.

\begin{cousinweb}
\emph{inoltrare} is built from \emph{in-} and \emph{oltre} \textquotedblleft{}beyond, further\textquotedblright{}, inherited from Latin \emph{ultrā} \textquotedblleft{}beyond\textquotedblright{} (Lewis \& Short): to send something further on.
\end{cousinweb}

\begin{grammarlens}[title={What you've built}]
One more word for telling, explaining and describing at B1.
\end{grammarlens}

\subsection*{Your turn}
\begin{itemize}
\raggedright
  \item \emph{Say it:} \emph{inoltrare}
  \item \emph{Say it:} \emph{inoltrare}, once more
  \item \emph{Say it:} a sentence of your own with \emph{inoltrare}
  \item \emph{Recall:} say the Italian for to publish, then the Italian for to comment, then say \emph{inoltrare} again
\end{itemize}

\begin{tcolorbox}[breakable,colback=teal!4,colframe=teal!35!black,title={Before you move on}]
What does inoltrare mean? (\textquotedblleft{}To forward (a message); to submit\textquotedblright{}.) Say it once more.
\end{tcolorbox}
\section[l'email]{l'email — the email}

\label{lesson:IT-C403-email}



\begin{quote}
Before the new one: say the Italian for to comment, then the Italian for to forward.
\end{quote}

\subsection*{You'll want to know: l'email}
\textbf{l'email} — \textquotedblleft{}the email\textquotedblright{}.

\begin{quote}
Se non rispondi all'email, ti chiamo.
\end{quote}

\begin{quote}
\emph{If you do not answer the email, I will call you.}
\end{quote}

Feminine: \emph{un'email, le email}. Also \emph{la mail} or \emph{la posta elettronica}.

\begin{cousinweb}
\emph{email} was borrowed from English \emph{e-mail}, short for \emph{electronic mail}; \emph{mail} goes back through Middle English \emph{male} to Anglo-Norman \emph{male} (EtymDB), usually explained as a word for a bag.
\end{cousinweb}

\begin{grammarlens}[title={What you've built}]
One more word for telling, explaining and describing at B1.
\end{grammarlens}

\subsection*{Your turn}
\begin{itemize}
\raggedright
  \item \emph{Say it:} \emph{l'email}
  \item \emph{Say it:} \emph{l'email}, once more
  \item \emph{Say it:} a sentence of your own with \emph{l'email}
  \item \emph{Recall:} say the Italian for to comment, then the Italian for to forward, then say \emph{l'email} again
\end{itemize}

\begin{tcolorbox}[breakable,colback=teal!4,colframe=teal!35!black,title={Before you move on}]
What does l'email mean? (\textquotedblleft{}The email\textquotedblright{}.) Say it once more.
\end{tcolorbox}
\section[la casella]{la casella — the inbox, the mailbox; the box (on a form)}

\label{lesson:IT-C403-casella}



\begin{quote}
Before the new one: say the Italian for to forward, then the Italian for the email.
\end{quote}

\subsection*{You'll want to know: la casella}
\textbf{la casella} — \textquotedblleft{}the inbox, the mailbox; the box (on a form)\textquotedblright{}.

\begin{quote}
Controlla la casella: ti ho scritto ieri.
\end{quote}

\begin{quote}
\emph{Check your inbox: I wrote to you yesterday.}
\end{quote}

\emph{La casella di posta} — the inbox. On a form, \emph{barrare la casella} — tick the box.

\begin{cousinweb}
\emph{casella} is a little \emph{casa} \textquotedblleft{}house\textquotedblright{}, inherited from Latin \emph{casa} \textquotedblleft{}a hut, a cottage\textquotedblright{} (Lewis \& Short): a small compartment, then a small box for letters.
\end{cousinweb}

\begin{grammarlens}[title={What you've built}]
One more word for telling, explaining and describing at B1.
\end{grammarlens}

\subsection*{Your turn}
\begin{itemize}
\raggedright
  \item \emph{Say it:} \emph{la casella}
  \item \emph{Say it:} \emph{la casella}, once more
  \item \emph{Say it:} a sentence of your own with \emph{la casella}
  \item \emph{Recall:} say the Italian for to forward, then the Italian for the email, then say \emph{la casella} again
\end{itemize}

\begin{tcolorbox}[breakable,colback=teal!4,colframe=teal!35!black,title={Before you move on}]
What does la casella mean? (\textquotedblleft{}The inbox, the mailbox; the box (on a form)\textquotedblright{}.) Say it once more.
\end{tcolorbox}
